\section{课程定位与问题定义}

本讲是 Berkeley Agentic AI MOOC F25 的安全专题开篇。Dawn Song 先强调一个背景：2025 年 Agentic AI 增长极快，能力边界持续外扩，系统能够感知环境、调用工具、执行外部动作，已经从 ``模型输出内容'' 走向 ``模型驱动行为''。这使得安全问题从传统 LLM 的对话安全，升级为系统级执行安全。

\begin{importantbox}{为什么 Agent 安全比纯 LLM 安全更难}
LLM 安全主要处理 ``输出是否有害''，Agent 安全还要处理 ``动作是否越权''、``工具是否被滥用''、``环境交互是否可控''。一旦 Agent 拥有写文件、发邮件、执行命令、访问网页和数据库权限，攻击后果就不再停留在文本层，而会落到真实世界系统。
\end{importantbox}

课程的主问题可归纳为三句：
\begin{enumerate}
    \item 如何定义 Agentic AI 的安全目标与安全边界？
    \item 如何系统化建模攻击面、威胁模型和风险评估流程？
    \item 如何建立可工程落地的纵深防御机制，而不是单点护栏？
\end{enumerate}

\begin{knowledgebox}{本讲议程}
讲座按 ``定义问题 $\rightarrow$ 拆解攻击 $\rightarrow$ 风险评估 $\rightarrow$ 防御设计 $\rightarrow$ 双用途风险'' 展开。它不是单一算法课，而是一门系统安全工程课。
\end{knowledgebox}

\subsection{本章小结}
本章的关键结论是：Agentic AI 安全的本质是系统安全问题，而不仅是模型内容安全问题。只讨论 prompt 或输出过滤无法覆盖真实风险面。

